%% kapa: exam=1 round=22 year=2011 period=2 form=A subject=appraisal_law subject_ko=부동산관계법규 seq=3 number=43
%% source: Q-Net 감정평가사 기출문제 (cst00309) artlSeq=5036790 file=extracted/1008/5036790/2011년 제22회 감정평가사 1차시험 기출문제/2교시.hwp
%% conversion: kordoc 4.24.51 -> tools/gen1.py (auto); review=proofread-gpt-astra; flags=none
\begin{question}{43}
국토의 계획 및 이용에 관한 법령상 도시관리계획에 관한 설명으로 옳지 않은 것은?
\begin{choices}
  \item 지구단위계획구역의 지정 또는 변경에 관한 계획은 도시관리계획에 속한다.
  \item 시ㆍ도지사는 도시관리계획을 결정하려면 관계 행정기관의 장과 미리 협의하여야 한다.
  \item 도시관리계획 입안권자는 도시관리계획을 조속히 입안하여야 할 필요가 있다고 인정되면 광역도시계획을 수립할 때에 도시관리계획을 입안할 수 있다.
  \item 도시관리계획 결정은 고시가 된 날부터 5일 후에 그 효력이 발생한다.
  \item 시가화조정구역의 지정에 관한 도시관리계획 결정 당시 이미 허가를 받아 사업이나 공사에 착수한 자는 별도의 신고 없이 그 사업이나 공사를 계속할 수 있다.
\end{choices}
\end{question}
\answer{5}
